% !TeX program = xelatex
\documentclass[11pt,a4paper]{article}

% ============================================================
% IAP Stage-1 Implementation Plan for Codex
% This document is designed to be provided directly to Codex.
% ============================================================

\usepackage[fontset=none]{ctex}
\setCJKmainfont[AutoFakeSlant=0.2]{Noto Serif CJK SC}
\setCJKsansfont{Noto Sans CJK SC}
\setCJKmonofont{Noto Sans Mono CJK SC}
\usepackage{amsmath,amssymb,amsfonts,mathtools,bm}
\usepackage{geometry}
\usepackage{booktabs,longtable,tabularx,array,multirow}
\usepackage{xcolor}
\usepackage{enumitem}
\usepackage{hyperref}
\usepackage{titlesec}
\usepackage{caption}
\usepackage{fancyhdr}
\usepackage{listings}
\usepackage{tcolorbox}
\tcbuselibrary{breakable,skins}

\geometry{margin=2.25cm}
\hypersetup{colorlinks=true, linkcolor=blue!55!black, urlcolor=blue!60!black, citecolor=blue!55!black}
\setlist[itemize]{leftmargin=1.5em,itemsep=2pt,topsep=2pt}
\setlist[enumerate]{leftmargin=1.8em,itemsep=2pt,topsep=2pt}
\newcolumntype{Y}{>{\raggedright\arraybackslash}X}
\newcolumntype{P}[1]{>{\raggedright\arraybackslash}p{#1}}
\renewcommand{\arraystretch}{1.18}

% -------------------- Colors --------------------
\definecolor{iapblue}{RGB}{0,80,150}
\definecolor{iapgreen}{RGB}{25,120,70}
\definecolor{iaporange}{RGB}{190,105,0}
\definecolor{iapred}{RGB}{180,40,40}
\definecolor{codebg}{RGB}{248,248,248}

\titleformat{\section}{\Large\bfseries\color{iapblue}}{\thesection}{0.7em}{}
\titleformat{\subsection}{\large\bfseries\color{iapblue!85!black}}{\thesubsection}{0.7em}{}
\titleformat{\subsubsection}{\normalsize\bfseries\color{iapblue!75!black}}{\thesubsubsection}{0.7em}{}

\setlength{\headheight}{14pt}
\pagestyle{fancy}
\fancyhf{}
\lhead{IAP Stage-1 v2.3 Implementation Plan}
\rhead{Codex Spec}
\cfoot{\thepage}

% -------------------- Commands --------------------
\newcommand{\PL}{\mathrm{PL}}
\newcommand{\AL}{\mathrm{AL}}
\newcommand{\IM}{\mathrm{IM}}
\newcommand{\HPL}{\mathrm{HPL}}
\newcommand{\VPL}{\mathrm{VPL}}
\newcommand{\HAL}{\mathrm{HAL}}
\newcommand{\VAL}{\mathrm{VAL}}
\newcommand{\Lam}{\Lambda}
\newcommand{\Sig}{\Sigma}
\newcommand{\Gm}{\mathbf{G}}
\newcommand{\Wm}{\mathbf{W}}
\newcommand{\Jm}{\mathbf{J}}
\newcommand{\Rm}{\mathbf{R}}
\newcommand{\evec}{\mathbf{e}}
\newcommand{\pvec}{\mathbf{p}}
\newcommand{\nm}{\mathbf{n}}
\newcommand{\xm}{\mathbf{x}}
\newcommand{\ym}{\mathbf{y}}
\newcommand{\eps}{\varepsilon}
\newcommand{\pos}{\mathrm{pos}}
\newcommand{\adv}{\mathrm{adv}}
\newcommand{\prior}{\mathrm{prior}}
\newcommand{\mon}{\mathrm{mon}}
\newcommand{\used}{\mathrm{used}}
\newcommand{\pred}{\mathrm{pred}}
\newcommand{\diag}{\mathrm{diag}}
\newcommand{\tr}{\mathrm{tr}}
\newcommand{\blkdiag}{\mathrm{blkdiag}}
\newcommand{\sat}{\mathrm{sat}}

\lstdefinestyle{codestyle}{
  basicstyle=\ttfamily\small,
  backgroundcolor=\color{codebg},
  frame=single,
  rulecolor=\color{black!20},
  breaklines=true,
  columns=fullflexible,
  keepspaces=true,
  showstringspaces=false,
  tabsize=2
}
\lstset{style=codestyle}

\newtcolorbox{codexbox}[1]{
  enhanced,
  breakable,
  colback=blue!2,
  colframe=iapblue!55!black,
  title=\textbf{#1},
  fonttitle=\bfseries,
  left=1mm,right=1mm,top=1mm,bottom=1mm
}

\newtcolorbox{warningbox}[1]{
  enhanced,
  breakable,
  colback=red!2,
  colframe=iapred!75!black,
  title=\textbf{#1},
  fonttitle=\bfseries,
  left=1mm,right=1mm,top=1mm,bottom=1mm
}

\newtcolorbox{acceptbox}[1]{
  enhanced,
  breakable,
  colback=green!2,
  colframe=iapgreen!65!black,
  title=\textbf{#1},
  fonttitle=\bfseries,
  left=1mm,right=1mm,top=1mm,bottom=1mm
}

% -------------------- Title --------------------
\title{\textbf{IAP Stage-1 v2.3 Implementation Plan for Codex}\\[3pt]
  \large Certified ARAIM Monitor $\parallel$ Advisory FIM-Fused Predictor\\[2pt]
\normalsize 完整公式、实现逻辑、接口规格、实现护栏、分阶段验收标准}
\author{Prepared for Codex-assisted implementation}
\date{May 2026}

\begin{document}
\maketitle
\vspace{-1.2em}

\begin{warningbox}{使用边界 / Read This First}
  本文档是给 Codex 单独执行开发用的 implementation plan，而不是论文方法章。Codex 必须先 inspect 当前 repository，再输出具体文件级 plan；未经确认，不得一次性重构全部系统。
  \begin{itemize}
    \item \textbf{必须保留 current certified ARAIM monitor 的 conservative max fusion。}
    \item \textbf{FIM-add 只允许进入 future/advisory predictor 与 planner path。}
    \item \textbf{每个 stage 单独 commit；每个 commit 只做一类改动。}
    \item \textbf{所有新公式必须有 log、config、fallback、unit test 或 launch-level sanity check。}
    \item \textbf{所有涉及坐标的实现必须显式标注 frame；不得混用 ENU/map/body/ECEF。}
    \item \textbf{所有旧 topic / message 必须在至少一个 stage 内保留兼容层，禁止一次性删除。}
  \end{itemize}
\end{warningbox}

\tableofcontents
\newpage


% ============================================================
\section{v2.3 增补摘要：本版相对 v2.2 的新增内容}
% ============================================================

本版在 v2.2 的基础上补齐 Codex 实现时最容易出错的公式边界、退化条件和执行护栏。新增重点如下：
\begin{enumerate}
  \item 明确 current LiDAR ARAIM 中 $\sigma_{ss}$ 的保守计算、数值检查和 fallback flag，避免 Codex 直接套用不适用于 LiDAR downdate 的 GNSS separation variance。
  \item 明确 advisory HPL/VPL 的 horizontal/vertical projection，尤其是 horizontal covariance 必须使用 $\lambda_{\max}(\Sigma_{xy})$，不能只取 $\max(\Sigma_{xx},\Sigma_{yy})$。
  \item 明确 GNSS FIM 的坐标系、clock bias 单位、Schur complement 退化条件和 PSD 检查。
  \item 明确 LiDAR $3\times3$ FIM 的最小可实现版本：先做 no-raycast local normal FIM，再逐步加入 raycast、incidence、dynamic 和 $6\times6$ pose FIM。
  \item 明确 A* front-end 与 B-spline back-end 对同一 $c_{PI}$ 的使用公式，防止两条 planner path 调出不同语义。
  \item 增加 stage invariants 与 Codex execution gate：Codex 必须先 inspect repo，并逐 stage 输出 file-level plan；未经批准不得直接改 Stage 2--4。
\end{enumerate}

\begin{acceptbox}{v2.3 的使用方式}
  这份文档已经可以作为 Codex 的 master implementation spec，但第一轮 Codex 任务仍然必须是 repository inspection，不是直接实现。每个 stage 都必须先输出实际文件、类、函数、topic、message、launch/test 命令和坐标系假设。
\end{acceptbox}

% ============================================================
\section{目标、范围与非目标}
% ============================================================

\subsection{最终目标}

本次整改目标是将当前 IAP Stage-1 系统从“可运行 demo”推进为“概念边界清晰、公式一致、可调试、可验证”的 integrity-aware planning 系统。核心结构为：
\[
  \boxed{
    \underbrace{\text{Certified ARAIM Monitor}}_{\text{current, safety-critical, conservative}}
    \quad \parallel \quad
    \underbrace{\text{Advisory FIM-Fused Predictor + Planner}}_{\text{future, optimization-friendly}}
  }
\]

\noindent 具体输出包括：
\begin{enumerate}
  \item 修复 LiDAR ARAIM PL 的非物理线性增长。
  \item 区分 current certified PL 与 future advisory PL proxy。
  \item 将 future GNSS/LiDAR prediction 从 two-channel PL/max logic 改为 position-FIM-level fusion。
  \item 将 PI function 改为只消费一个统一的 advisory predicted PL / integrity margin。
  \item 将 ESDF、AL、PL、IM、PI cost 合并为 Unified Risk Grid (URG) 的多层查询结构。
  \item 建立 stage-by-stage 的 log、test、验收标准和性能 profiling。
\end{enumerate}

\subsection{非目标}

本次整改\textbf{不做}以下事情：
\begin{itemize}
  \item 不将 SLAM/FGO 与 planner 做完全联合优化。
  \item 不把 FIM-add 用于 certified/current ARAIM monitor。
  \item 不声称 future predicted PL proxy 是严格 certified ARAIM PL。
  \item 不一次性切换到完整 $6\times 6$ LiDAR pose-level FIM；第一版只要求 $3\times 3$ translational FIM。
  \item 不默认开启 balance/fault-robust ablation 项；第一版主线只启用 PI main cost。
\end{itemize}

\subsection{Codex 执行方式}

\begin{codexbox}{Codex 必须遵守的工作流}
  \begin{enumerate}
    \item Inspect repository and identify actual files/classes/topics related to ARAIM, PL predictor, PI cost, EGO planner integration, ESDF/map, coordinate frames, and logging.
    \item Produce a file-level implementation plan for the current stage only.
    \item Fill the repository-specific command template: build command, unit test command, launch sanity command, log extraction command, and plotting/debug command.
    \item Wait for confirmation before modifying code.
    \item Implement one stage per commit; do not mix Stage 0 with Stage 2/3/4.
    \item After each stage, run available build/tests/launch sanity checks and summarize exact changed files, configs, logs, coordinate-frame assumptions, and known limitations.
  \end{enumerate}
\end{codexbox}

% ============================================================
\section{系统主架构：双路径分离}
% ============================================================

\subsection{Current certified monitor path}

Current monitor path 使用真实当前时刻测量、当前 FGO snapshot、真实 residual 或 subset solution separation，输出 current GNSS PL、current LiDAR PL 与 alarm flag。其融合保持保守：
\[
  \boxed{
    \PL^{\mon}_q = \max\left(\PL^G_q,\ \PL^L_q\right),
    \qquad q\in\{H,V,E,N,U\}.
  }
\]

\noindent 该路径的职责是报警与安全监督，不承担 planner cost field 的高频查询。Codex 不得将该路径改为 FIM-add。

\subsection{Future advisory predictor path}

Future predictor path 对候选位置 $\pvec$、未来时间 $t+\tau$、局部地图、卫星几何、当前 covariance prior 进行预测，输出 advisory predicted PL proxy 与 integrity margin field：
\[
  \Lam^{\adv}_{\pos}(\pvec,t)
  =
  \Lam^{\prior}_{\pos}(t)
  +
  \Lam^{G}_{\pos}(\pvec,t)
  +
  \Lam^{L}_{\pos}(\pvec,t).
\]
\[
  \Sig^{\adv}_{\pos}(\pvec,t)
  =
  \left(\Lam^{\adv}_{\pos}(\pvec,t)+\lambda_{\mathrm{reg}}I\right)^{-1}.
\]
\[
  \PL^{\adv}(\pvec,t)
  \rightarrow
  \IM(\pvec,t)=\AL(\pvec,t)-\PL^{\adv}(\pvec,t)
  \rightarrow
  c_{PI}(\pvec,t).
\]

\subsection{命名规则}

\begin{center}
  \small
  \begin{tabularx}{\linewidth}{l l Y}
    \toprule
    \textbf{旧名称} & \textbf{新名称} & \textbf{说明} \\
    \midrule
    GNSS PL & GNSS Certified PL & current monitor 输出，可用于 alarm。\\
    LiDAR PL & LiDAR Certified PL & current VGICP-block ARAIM 输出，可用于 alarm。\\
    Predicted GNSS PL & GNSS Advisory PL Proxy / GII & future geometry-driven, non-certified。\\
    Predicted LiDAR PL & LiDAR Advisory PL Proxy / LOI & future FIM-driven, non-certified。\\
    Predicted Fused PL & Advisory Predicted PL & future FIM-add 后导出的米级 proxy。\\
    Integrity Margin & Advisory Integrity Margin & $\IM=\AL-\PL^{\adv}$。\\
    \bottomrule
  \end{tabularx}
\end{center}


% ============================================================
\section{Codex 实现护栏：必须先读}
% ============================================================

\subsection{Stage-level hard Do / Do-not rules}

\begin{warningbox}{Global hard constraints}
Codex 必须把本节当作比公式更高优先级的实现约束。若当前代码结构与本 spec 冲突，Codex 先报告冲突并提出最小兼容方案，不得直接大改架构。
\end{warningbox}

\begin{center}
\small
\begin{tabularx}{\linewidth}{P{0.15\linewidth}Y Y}
  \toprule
  \textbf{Stage} & \textbf{Allowed} & \textbf{Forbidden} \\
  \midrule
  Stage 0 & bounded age, fixed target window, component logs, config defaults & changing planner behavior; changing current monitor fusion; introducing FIM-add; renaming topics unless unavoidable \\
  Stage 1 & comments, log labels, documentation, non-breaking aliases & numerical behavior changes; deleting old fields/topics; changing alarm logic \\
  Stage 2 & advisory GNSS/LiDAR FIM, FIM-add in predictor only, old mode kept for A/B & modifying certified/current ARAIM monitor; changing alarm logic; requiring raycast in first implementation; introducing $6\times6$ LiDAR FIM; changing planner cost \\
  Stage 3 & PI main cost using $\PL^{\adv}$, $\AL$, $\IM$; optional terms behind disabled flags & enabling ratio/balance/fault-robust by default; silently tuning planner weights; changing ESDF/collision cost; consuming current certified PL as planner field \\
  Stage 4 & URG compatibility layer, new query API, stale unknown-risk handling & deleting old topics before validation; replacing A* and B-spline query paths in one patch; stale-to-zero fallback \\
  Stage 5 & cache, sparse active voxels, parallel refresh, profiling & algorithmic changes that alter Stage 2/3 semantics without A/B switch \\
  \bottomrule
\end{tabularx}
\end{center}

\subsection{Coordinate-frame conventions}

所有 advisory grid query points 使用 world/map frame $W$：
\[
  \pvec_W=[x_W,y_W,z_W]^\top.
\]
Codex 必须在代码注释和日志里标注以下 frame 假设：
\begin{itemize}
  \item A* and B-spline sample/control points passed to \texttt{queryRisk()} are in $W$.
  \item ESDF/occupancy distance query is in $W$.
  \item LiDAR voxel centers $c_m$ and normals $\nm_m$ used by advisory FIM must be in $W$.
  \item GNSS satellite positions must be converted to the same local navigation frame as $W$ before line-of-sight computation. If the repository currently uses ENU/ECEF conversion elsewhere, reuse that converter; do not invent a new frame transform.
  \item If any frame transform is unavailable, Codex must stop and report the missing transform instead of implementing a silent approximation.
\end{itemize}

\subsection{Fallback and compatibility policy}

\begin{itemize}
  \item Existing topics such as \texttt{/iap/integrity\_cost\_field} and \texttt{/iap/integrity\_front\_cost\_field} must remain available until Stage 4 is validated.
  \item Before URG is fully integrated, old planner-consumable cost-field topics may be generated from the new $\PL^{\adv}$ and $c_{PI}$ pipeline through a compatibility publisher.
  \item Missing or degenerate source contributions must set flags and produce a safe fallback; they must never produce NaN/Inf.
  \item Stale advisory field must be treated as unknown risk, not zero risk.
\end{itemize}

\subsection{Repository-inspection output template}

Codex 的第一轮回复必须包含以下模板内容，不得直接改代码：
\begin{lstlisting}
Stage to inspect: <0/1/2/3/4/5>
Relevant files/classes/functions:
- ARAIM monitor:
- LiDAR block risk score:
- Target keyframe selection:
- PL composition/fusion:
- Predictor/grid field:
- PI cost adapter:
- Planner A* integration:
- B-spline backend integration:
- ESDF/occupancy map:
- Message/topic definitions:
- Config files:
- Logging utilities:

Coordinate-frame assumptions found in code:
- Planner frame:
- ESDF/map frame:
- LiDAR voxel/map frame:
- GNSS satellite/local frame:
- Existing transform utilities:

Available commands discovered:
- Build:
- Unit tests:
- Launch sanity check:
- Log extraction:
- Plot/debug visualization:

Minimal patch plan:
Risks/unknowns:
Questions before implementation:
\end{lstlisting}

% ============================================================
\section{核心数学定义与公式包}
% ============================================================

\subsection{当前 ARAIM monitor 的 LiDAR PL 公式}

对当前 FGO 线性化点 $\hat{x}_0$，每个 VGICP factor block $F_{j,\ell}$ 贡献：
\[
  \Delta \Lam_{j,\ell}=J_{j,\ell}^\top R_{j,\ell}^{-1}J_{j,\ell},
  \qquad
  \Delta g_{j,\ell}=J_{j,\ell}^\top R_{j,\ell}^{-1}r_{j,\ell}.
\]

对 fault hypothesis $H_f$，snapshot downdate 为：
\[
  \Lam_f=\Lam_0-\sum_{F_{j,\ell}\in B(H_f)}\Delta \Lam_{j,\ell},
  \qquad
  g_f=g_0-\sum_{F_{j,\ell}\in B(H_f)}\Delta g_{j,\ell}.
\]
\[
  \delta x_f=\Lam_f^{-1}g_f,
  \qquad
  \Sig_f=\Lam_f^{-1},
  \qquad
  d_f=\hat{x}_0\ominus \hat{x}_f.
\]

LiDAR current PL 保持 solution-separation 形式：
\[
  \boxed{
    \PL^L_{q,f}
    =
    |e_q^\top d_f|
    +K_{\mathrm{fa},f}\sigma_{\mathrm{ss},q,f}
    +K_{\mathrm{md},f}\sigma_{q,f}
    +b_{q,f},
    \qquad
    \PL^L_q=\max_f\PL^L_{q,f}.
  }
\]


\subsection{Solution-separation variance 的实现级定义与 fallback}

Current LiDAR ARAIM 使用 snapshot downdate，而不是单历元 WLS subset solver。因此 Codex 不得无条件套用 GNSS ARAIM 中的投影矩阵公式。实现规则如下。

若当前代码已经有严格的 solution-separation covariance routine，例如
\[
  \Sigma_{ss,f}=(S_0-S_f)C_y(S_0-S_f)^\top,
\]
则保持现有实现，但必须加数值检查和日志。若当前只暴露 full/subset covariance，则使用以下保守近似：
\[
  \sigma^2_{ss,q,f,raw}=e_q^\top(\Sigma_f-\Sigma_0)e_q.
\]
实现时必须检查：
\[
  \mathrm{finite}(\sigma^2_{ss,q,f,raw})\quad\text{and}\quad \sigma^2_{ss,q,f,raw}>\sigma^2_{ss,min}.
\]
若检查通过：
\[
  \boxed{\sigma_{ss,q,f}=\sqrt{\max(\sigma^2_{ss,q,f,raw},\sigma^2_{ss,min})}.}
\]
若检查失败、$\Sigma_f-\Sigma_0$ 非正、或 subset covariance 数值不可信，则 fallback 为：
\[
  \boxed{\sigma_{ss,q,f}=\sqrt{\max(e_q^\top\Sigma_f e_q,\sigma^2_{ss,min})}.}
\]
并设置：
\[
  \texttt{flags}\leftarrow\texttt{SS\_VARIANCE\_FALLBACK}\cup\texttt{FIM\_REGULARIZED}.
\]

\begin{warningbox}{Why this fallback is required}
  LiDAR factor-block downdate 不保证与传统 GNSS WLS informative subset 关系完全一致。若 Codex 发现 $\Sigma_f-\Sigma_0$ 的投影为负，不得取绝对值悄悄修复；必须 fallback、打 flag、记录 worst hypothesis 和对应 eigenvalue。
\end{warningbox}

推荐实现伪代码：
\begin{lstlisting}[language=C++]
double computeSigmaSs(const Matrix3d& Sigma0,
                      const Matrix3d& Sigmaf,
                      const Vector3d& eq,
                      double sigma_ss_min,
                      uint32_t* flags) {
  const double raw = (eq.transpose() * (Sigmaf - Sigma0) * eq)(0,0);
  if (std::isfinite(raw) && raw > sigma_ss_min * sigma_ss_min) {
    return std::sqrt(std::max(raw, sigma_ss_min * sigma_ss_min));
  }
  const double fb = (eq.transpose() * Sigmaf * eq)(0,0);
  *flags |= SS_VARIANCE_FALLBACK | FIM_REGULARIZED;
  return std::sqrt(std::max(fb, sigma_ss_min * sigma_ss_min));
}
\end{lstlisting}

\subsection{LiDAR block risk score 稳定化}

旧式风险分数包含无界 age 项：
\[
  \gamma_{j,\ell}^{old}
  =w_r\gamma_{\mathrm{rmse}}
  +w_i\gamma_{\mathrm{inlier}}
  +w_c\gamma_{\kappa}
  +w_a\frac{\mathrm{Age}_j}{A_{ref}}.
\]

新式 age 项必须有界：
\[
  \boxed{
    \gamma_{\mathrm{age}}(A)=1-e^{-A/\tau_A}
    \quad \text{or}\quad
    \gamma_{\mathrm{age}}(A)=\min(A/A_{ref},\gamma_{\max}).
  }
\]

新式 block score：
\[
\begin{aligned}
  \gamma_{j,\ell}
  ={}&
  w_r\,\mathrm{clip}\left(\frac{\mathrm{RMSE}_{j,\ell}}{\mathrm{RMSE}_{ref}},0,\gamma_r^{\max}\right)
  +w_i\,(1-\mathrm{InlierFrac}_{j,\ell}) \\
  &+w_c\,\mathrm{clip}\left(\frac{\log\kappa_{j,\ell}}{\log\kappa_{ref}},0,\gamma_c^{\max}\right)
  +w_a\gamma_{\mathrm{age}}(\mathrm{Age}_j).
\end{aligned}
\]

Bias overbound 仍保留工程形式，但必须 log 单项：
\[
  b_{q,f}=\alpha_q\max_{F_{j,\ell}\in B(H_f)}\gamma_{j,\ell}.
\]

\subsection{Target keyframe 固定窗口}

为消除 $\max$ over growing set 带来的集合膨胀，当前 LiDAR ARAIM 与 LiDAR predictor 中使用的 target keyframe 集合必须固定上限：
\[
  \boxed{|T_t|\le K_T.}
\]

默认：
\[
  K_T=10.
\]
Target selection 建议排序优先级：
\begin{enumerate}
  \item spatial proximity to current pose；
  \item temporal freshness；
  \item overlap / inlier quality；
  \item avoid duplicate near-identical keyframes。
\end{enumerate}

\subsection{Prior position FIM from covariance}

当前 advisory predictor 只使用 position-domain FIM。若 estimator 提供 position covariance $\Sig^{prior}_{pos}$：
\[
  \boxed{
    \Lam^{\prior}_{pos}
    =
    \left(\eta_\Sigma\Sig^{prior}_{pos}+\sigma_{floor}^2I\right)^{-1}.
  }
\]

其中 $\eta_\Sigma\ge 1$ 是 covariance inflation，$\sigma_{floor}$ 防止 iSAM covariance 过小或奇异。默认：
\[
  \eta_\Sigma=2.0,
  \qquad
  \sigma_{floor}=0.05\ \mathrm{m}.
\]

若只提供 $6\times6$ pose covariance，取 position block：
\[
  \Sig^{prior}_{pos}=\Sig^{pose}_{1:3,1:3}.
\]
若 covariance 缺失、含 NaN/Inf 或非正定，则使用配置项 \texttt{default\_prior\_sigma\_xyz\_m} 构造 diagonal fallback covariance，并设置 \texttt{FIM\_REGULARIZED}。

\subsection{GNSS future advisory FIM}

对候选位置 $\pvec$ 与时间 $t$，可见卫星集合为 $\mathcal{S}(\pvec,t)$。第 $s$ 颗卫星的 line-of-sight unit vector：
\[
  \ell_s(\pvec,t)=\frac{p_s^{sat}(t)-\pvec}{\|p_s^{sat}(t)-\pvec\|}.
\]

GNSS geometry matrix 含接收机 clock bias：
\[
  G(\pvec,t)=
  \begin{bmatrix}
    -\ell_1^\top & 1\\
    -\ell_2^\top & 1\\
    \vdots & \vdots\\
    -\ell_N^\top & 1
  \end{bmatrix}
  \in\mathbb{R}^{N\times 4}.
\]

预测测量方差：
\[
  \sigma^2_{eff,s}
  =
  \sigma^2_{URE,s}
  +\sigma^2_{tropo,s}
  +\sigma^2_{iono,s}
  +\sigma^2_{mp,s}(\pvec,el_s)
  +\sigma^2_{NLOS,s}(\pvec)
  +\sigma^2_{canopy,s}(\pvec).
\]
\[
  W(\pvec,t)=\diag(\sigma^{-2}_{eff,1},\dots,\sigma^{-2}_{eff,N}).
\]

Full GNSS Hessian：
\[
  H_G=G^\top WG=
  \begin{bmatrix}
    H_{pp} & H_{pc}\\
    H_{cp} & H_{cc}
  \end{bmatrix},
  \qquad
  H_{pp}\in\mathbb{R}^{3\times3},\ H_{cc}\in\mathbb{R}^{1\times1}.
\]

消去 clock bias 后的 position-domain GNSS FIM：
\[
  \boxed{
    \Lam^G_{pos}(\pvec,t)=H_{pp}-H_{pc}(H_{cc}+\lambda_c)^{-1}H_{cp}.
  }
\]

若 $N<4$ 或 $H_{cc}$ 退化，则设置：
\[
  \Lam^G_{pos}=0,
  \qquad
  \texttt{flags} \leftarrow \texttt{GNSS\_DEGENERATE}.
\]
若 satellite variance 缺失或非正，使用 \texttt{sigma\_gnss\_default\_m}；若 satellite position 不可用，跳过该卫星并在 refresh log 中记录 skipped count。


\subsection{GNSS FIM 的坐标系、clock 单位与 PSD 检查}

Codex 必须在实现 GNSS advisory FIM 前确认当前代码的坐标链。本文公式
\[
  G_s=[-\ell_s^\top\quad 1]
\]
默认 pseudorange residual 单位为 meter，clock variable 为 meter-equivalent receiver clock bias：
\[
  b^{clk}=c\delta t^{clk}\quad [\mathrm{m}].
\]
若当前 estimator 内部 clock state 使用 second，则 geometry row 必须改为：
\[
  G_s=[-\ell_s^\top\quad c],
\]
或者先把 clock state 转为 meter-equivalent。两种写法不得混用。

Satellite position 必须先转换到 planner/advisory grid 使用的 local frame：
\[
  p^{sat}_{local}=T_{ECEF\rightarrow local}\,p^{sat}_{ECEF},
  \qquad
  \pvec\in\mathrm{same\ local\ frame}.
\]
若 transform unavailable 或 frame id 不一致，则该 refresh 不应生成 GNSS FIM：
\[
  \Lam^G_{pos}=0,
  \qquad
  \texttt{flags}\leftarrow\texttt{FRAME\_MISMATCH}\cup\texttt{GNSS\_DEGENERATE}.
\]

Schur complement 后必须做对称化与 PSD 检查：
\[
  \Lam^G_{pos}\leftarrow\frac{1}{2}\left(\Lam^G_{pos}+(\Lam^G_{pos})^\top\right).
\]
若
\[
  \lambda_{min}(\Lam^G_{pos})<-\epsilon_{PSD},
\]
则设置 \texttt{GNSS\_DEGENERATE}，并将 GNSS contribution 置零或 eigenvalue-clamp 到非负，具体由配置项 \texttt{gnss\_fim\_psd\_policy} 控制。默认第一版使用 \texttt{zero\_with\_flag}，避免用错误几何误导 planner。

\subsection{LiDAR future advisory FIM：第一版 \texorpdfstring{$3\times3$}{3x3} translational approximation}

候选点 $\pvec$ 周围潜在可见 voxel/normal 集合：
\[
  V(\pvec)=\{m:\ \|c_m-\pvec\|\le R_L,\ \mathrm{visible}(\pvec,c_m)=1,\ \mathrm{validNormal}(m)=1\}.
\]

第一版 LiDAR FIM：
\[
  \boxed{
    \Lam^L_{pos}(\pvec)
    =
    \sum_{m\in V(\pvec)}
    \pi_m(\pvec)\,w_m(\pvec)\,\nm_m\nm_m^\top.
  }
\]

其中：
\[
  w_m(\pvec)=\frac{1}{\sigma_m^2(\pvec)}.
\]

推荐的 visibility/overlap 权重：
\[
  \pi_m(\pvec)=
  \pi_{ray}(\pvec,c_m)\cdot
  \pi_{range}(\|c_m-\pvec\|)\cdot
  \pi_{age}(A_m)\cdot
  \pi_{dynamic}(m).
\]

默认各项：
\[
  \pi_{ray}=
  \begin{cases}1,&\text{ray not occupied}\\ 0,&\text{ray blocked}
  \end{cases}
\]
\[
  \pi_{range}(r)=\exp\left(-\frac{r^2}{2r_0^2}\right),
  \qquad
  \pi_{age}(A)=\exp(-A/\tau_{map}),
  \qquad
  \pi_{dynamic}=1-P_{dyn}(m).
\]

测量方差模型：
\[
  \sigma_m^2(\pvec)
  =
  \sigma_0^2
  +k_r r_m^2
  +k_i\frac{1}{\max(|\nm_m^\top v_{ray}|,\eps_i)}
  +k_n\sigma^2_{normal,m}
  +k_a A_m
  +k_d P_{dyn}(m).
\]

\begin{warningbox}{LiDAR FIM 的实现边界}
  第一版只要求 position-domain $3\times3$ translational FIM。Codex 不得在第一版把它写成 certified LiDAR ARAIM，也不得替代 current LiDAR ARAIM monitor。论文/注释中必须标注：\textit{lightweight translational approximation for advisory planning only}。

  第一版不要求 raycast、dynamic probability、GMM normals 或完整 $6\times6$ pose FIM。若 voxel normal statistics 不可用，必须使用 \texttt{lidar\_fallback\_mode}，而不是临时重构地图模块。
\end{warningbox}


\subsection{LiDAR FIM 最小实现阶梯}

为了避免 Codex 一次性重构地图系统，LiDAR advisory FIM 必须按以下阶梯实现。

\paragraph{Version A: no-raycast local normal FIM.}
第一版只使用候选点邻域内已有 voxel normal statistics：
\[
  V_A(\pvec)=\{m:\|c_m-\pvec\|<R_L,\ \mathrm{validNormal}(m)=1\},
  \qquad \pi_{ray}=1.
\]
\[
  \Lam^L_A(\pvec)=\sum_{m\in V_A(\pvec)}
  \exp\left(-\frac{\|c_m-\pvec\|^2}{2r_0^2}\right)
  \frac{1}{\sigma_m^2}\,n_m n_m^\top.
\]
若 normal histogram 不存在，则 Codex 必须复用当前 scalar advisory predictor 或返回 zero-with-flag，而不是临时实现新的 map extraction pipeline。

\paragraph{Version B: occupancy raycast visibility.}
在 Version A 稳定后加入 raycast：
\[
  \pi_{ray}(\pvec,c_m)=\mathbb{I}\{\mathrm{ray}(\pvec,c_m)\ \mathrm{not\ occupied}\}.
\]
该版本必须有单独配置 \texttt{lidar\_raycast\_enabled}，默认保持 false，直到 corridor test 与 latency test 通过。

\paragraph{Version C: incidence, dynamic, map-age weights.}
加入入射角、dynamic probability、map age：
\[
  \sigma_m^2=\sigma_0^2+k_r r_m^2+k_i/\max(|n_m^\top v_{ray}|,\epsilon_i)+k_a A_m+k_dP_{dyn}(m).
\]

\paragraph{Version D: future work, $6\times6$ pose FIM.}
完整 scan-matching FIM 应为
\[
  \Lam_L^{pose}=\sum_m J_m^\top R_m^{-1}J_m\in\mathbb{R}^{6\times6},
  \qquad
  \Sigma_p=[(\Lam_L^{pose})^{-1}]_{p,p}.
\]
Stage 2 禁止实现 Version D；只允许在注释和 future-work 中保留接口占位。

\subsection{Advisory FIM-add fusion 与 PL 导出}

统一 position FIM：
\[
  \boxed{
    \Lam^{\adv}_{pos}(\pvec,t)
    =\Lam^{\prior}_{pos}(t)
    +\Lam^G_{pos}(\pvec,t)
    +\Lam^L_{pos}(\pvec,t).
  }
\]

正则化逆：
\[
  \boxed{
    \Sig^{\adv}_{pos}(\pvec,t)
    =\left(\Lam^{\adv}_{pos}(\pvec,t)+\lambda_{reg}I\right)^{-1}.
  }
\]

水平 covariance：
\[
  \Sig^{\adv}_{xy}=\begin{bmatrix}
  \Sig^{\adv}_{xx} & \Sig^{\adv}_{xy}\\
  \Sig^{\adv}_{yx} & \Sig^{\adv}_{yy}
\end{bmatrix}.
\]

HPL proxy：
\[
\boxed{
  \HPL^{\adv}(\pvec,t)
  =
  K_H\sqrt{\lambda_{\max}(\Sig^{\adv}_{xy})}
  +b_H^{\pred}(\pvec,t)+s_H^{\pred}(\pvec,t).
}
\]

VPL proxy：
\[
\boxed{
  \VPL^{\adv}(\pvec,t)
  =
  K_V\sqrt{\Sig^{\adv}_{zz}}
  +b_V^{\pred}(\pvec,t)+s_V^{\pred}(\pvec,t).
}
\]

Unified scalar PL for planner：
\[
\boxed{
  \PL^{\adv}(\pvec,t)=\max\left(\HPL^{\adv}(\pvec,t),\ \VPL^{\adv}(\pvec,t)\right)
}
\]
或者 planner 可分别使用 horizontal/vertical constraints；第一版为了稳定，采用 max scalar。

默认第一版 bias/separation upper bound：
\[
b_q^{\pred}=b_{q,0}+b_{q,G}^{\pred}+b_{q,L}^{\pred},
\qquad
s_q^{\pred}=n_{ss}\sqrt{e_q^\top\Sig^{\adv}_{pos}e_q}.
\]
其中第一版可设置 $b_{q,G}^{\pred}, b_{q,L}^{\pred}$ 为可配置常数或由 quality score 线性映射，但必须 log。

\subsection{Alert Limit field}

水平 alert limit：
\[
\HAL(\pvec)=\gamma_H\max\left(d_{obs}(\pvec)-r_{drone}-r_{buffer},\ 0\right).
\]

垂直 alert limit：
\[
\VAL(\pvec)=\gamma_V\max\left(\min(z-z_{min},\ z_{max}-z),\ 0\right).
\]

统一 alert limit：
\[
\boxed{
  \AL(\pvec)=\min(\HAL(\pvec),\VAL(\pvec)).
}
\]

Integrity margin：
\[
\boxed{
  \IM(\pvec,t)=\AL(\pvec)-\PL^{\adv}(\pvec,t).
}
\]

\subsection{PI function}

主项默认只依赖统一 PL/AL：
\[
\boxed{
  c^{main}_{PI}(\pvec,t)
  =
  \lambda_\pi\left[\PL^{\adv}(\pvec,t)-\AL(\pvec)+m\right]_+^2.
}
\]

可选 ratio term：
\[
c^{ratio}_{PI}(\pvec,t)
=\mu_\pi\left(\frac{\PL^{\adv}(\pvec,t)}{\AL(\pvec)+\eps_{AL}}\right)^2.
\]

总代价第一版：
\[
\boxed{
  c_{PI}^{base}=c_{PI}^{main}+\mathbb{I}_{ratio}\,c_{PI}^{ratio}.
}
\]

Ablation 项默认关闭：
\[
r_G=\frac{\tr(\Lam_G)}{\tr(\Lam_G+\Lam_L)+\eps},
\qquad
c_{balance}=w_b[\max(r_G,1-r_G)-r_{max}]_+^2.
\]
\[
c_{fault}=w_f\left([\PL^{\adv}_{\setminus G}-\AL]_+^2+[\PL^{\adv}_{\setminus L}-\AL]_+^2\right).
\]


\subsection{Planner cost interface：A* 与 B-spline 的精确定义}

Stage 3 后，front-end A* 与 back-end B-spline optimizer 必须查询同一个 $c_{PI}$ 语义。Codex 不得分别写两套 GNSS/LiDAR source-specific cost。

A* edge cost 使用乘性 inflation：
\[
  \boxed{
  c_{edge}(p_i,p_j)
  =\|p_i-p_j\|
  \left(1+\lambda_{obs}c_{obs}(p_j)+\lambda_{PI}^{front}c_{PI}^{used}(p_j)\right).
  }
\]
其中 $c_{obs}$ 来自现有 ESDF/occupancy cost，Codex 不得在 Stage 3 自动调它的权重。

B-spline back-end 使用 sample/control-point penalty：
\[
  \boxed{
  J_{back}=J_{smooth}+J_{collision}+J_{dyn}
  +\lambda_{PI}^{back}\sum_{k} c_{PI}^{used}(Q_k).
  }
\]
其中 $Q_k$ 是当前后端已有的轨迹采样点或控制点查询位置。第一版梯度允许使用 finite difference 或现有 grid gradient：
\[
  \nabla c_{PI}(p)\approx
  \left[
  \frac{c_{PI}(p+h e_x)-c_{PI}(p-h e_x)}{2h},
  \frac{c_{PI}(p+h e_y)-c_{PI}(p-h e_y)}{2h},
  \frac{c_{PI}(p+h e_z)-c_{PI}(p-h e_z)}{2h}
  \right]^\top.
\]
若 finite-difference query 命中 stale/out-of-range voxel，则必须使用 query API 返回的 flags 决定 fallback，不得返回零梯度后继续认为该方向安全。

\begin{acceptbox}{Planner interface acceptance}
  PI disabled 时，A* 和 B-spline 应退回原始 EGO 行为；PI enabled 且 field fresh 时，front-end 与 back-end 的 $c_{PI}$ 数值应在同一 query point 上一致，允许仅因插值方式产生小误差。
\end{acceptbox}

\subsection{Unified Risk Grid 与 stale handling}

每个 voxel 存储：
\[
\boxed{
  URG(v)=\{d_{ESDF},\rho_{occ},\AL,\HPL^{\adv},\VPL^{\adv},\PL^{\adv},\IM,c_{PI},age,flags\}.
}
\]

Stale field 不允许 fallback 到 zero risk。使用 unknown penalty：
\[
\boxed{
  c_{PI}^{\used}(\pvec,t)=c_{PI}(\pvec,t)+c_{unknown}\left(1-e^{-\Delta t/\tau_{stale}}\right).
}
\]

若 $\Delta t>t_{max\_age}$：
\[
\texttt{flags}\leftarrow\texttt{STALE\_PL}\cup\texttt{UNKNOWN\_RISK},
\]
并触发 refresh / slow-down / advisory PI disable according to config。

% ============================================================
\section{配置参数默认值}
% ============================================================

\begin{center}
\small
\begin{longtable}{@{}P{0.34\linewidth}P{0.18\linewidth}P{0.42\linewidth}@{}}
  \toprule
  \textbf{参数} & \textbf{默认值} & \textbf{说明} \\
  \midrule
  \texttt{certified\_monitor\_fusion} & \texttt{max} & 禁止改为 FIM-add。\\
  \texttt{advisory\_fusion\_mode} & \texttt{fim\_add} & Predictor/planner path 默认。\\
  \texttt{enable\_advisory\_fim\_fusion} & \texttt{true} & Stage 2 完成后启用。\\
  \texttt{target\_window\_K} & 10 & LiDAR target keyframe 固定窗口。\\
  \texttt{age\_model} & \texttt{exp\_saturating} & $1-e^{-A/\tau_A}$。\\
  \texttt{age\_tau\_s} & 30.0 & LiDAR block age 饱和时间常数。\\
  \texttt{gamma\_age\_max} & 1.0 & 若使用 cap 模型。\\
  \texttt{covariance\_inflation\_eta} & 2.0 & prior covariance inflation。\\
  \texttt{sigma\_floor\_m} & 0.05 & 防止 covariance 过小。\\
  \texttt{default\_prior\_sigma\_xyz\_m} & 2.0,2.0,1.0 & covariance missing 时的 fallback prior。\\
  \texttt{fim\_regularization} & $10^{-6}$ & FIM inverse regularization。\\
  \texttt{sigma\_ss\_min\_m} & 0.02 & solution-separation std floor。\\
  \texttt{gnss\_fim\_psd\_policy} & \texttt{zero\_with\_flag} & GNSS Schur complement 非 PSD 时的默认策略。\\
  \texttt{lambda\_clock\_schur} & $10^{-6}$ & GNSS clock Schur complement regularization。\\
  \texttt{gnss\_clock\_unit} & \texttt{meter} & 若代码使用 seconds，必须显式乘 $c$ 并打日志。\\
  \texttt{sigma\_gnss\_default\_m} & 5.0 & missing GNSS variance fallback。\\
  \texttt{lidar\_fim\_radius\_m} & 12.0 & 候选点周围 voxel 搜索半径。\\
  \texttt{lidar\_range\_r0\_m} & 8.0 & range weight 衰减尺度。\\
  \texttt{lidar\_raycast\_enabled} & \texttt{false} first & 初版可 false，后续打开。\\
  \texttt{lidar\_min\_valid\_normals} & 6 & LiDAR FIM 最小有效 normal 数。\\
  \texttt{lidar\_fim\_min\_eig} & $10^{-5}$ & LiDAR FIM degeneracy threshold。\\
  \texttt{lidar\_fim\_max\_cond} & $10^6$ & LiDAR FIM condition threshold。\\
  \texttt{lidar\_fallback\_mode} & \texttt{zero\_with\_flag} & normal 缺失时 fallback；可选 old\_scalar。\\
  \texttt{normal\_hist\_az\_bins} & 8 & voxel normal histogram azimuth bins。\\
  \texttt{normal\_hist\_el\_bins} & 8 & voxel normal histogram elevation bins。\\
  \texttt{K\_H} & 5.0 & advisory HPL multiplier，后续 calibration。\\
  \texttt{K\_V} & 5.0 & advisory VPL multiplier，后续 calibration。\\
  \texttt{pi\_margin\_m} & 0.5 & PI hinge safety margin。\\
  \texttt{pi\_lambda} & 1.0 & PI hinge weight。\\
  \texttt{pi\_lambda\_front} & 1.0 & A* front-end PI edge inflation weight。\\
  \texttt{pi\_lambda\_back} & 1.0 & B-spline back-end PI penalty weight。\\
  \texttt{pi\_use\_ratio} & \texttt{false} first & ratio 项默认关闭。\\
  \texttt{pi\_mu\_ratio} & 0.1 & ratio 项权重。\\
  \texttt{pi\_use\_balance} & \texttt{false} & ablation only。\\
  \texttt{pi\_use\_fault\_robust} & \texttt{false} & ablation only。\\
  \texttt{pi\_max\_cost} & 100.0 & cost clipping。\\
  \texttt{unknown\_penalty} & 10.0 & stale field unknown risk penalty。\\
  \texttt{stale\_tau\_s} & 2.0 & unknown penalty time constant。\\
  \texttt{max\_field\_age\_s} & 2.0 & 超过后标记 stale。\\
  \texttt{urg\_resolution\_m} & 1.0 & 初版 grid resolution。\\
  \texttt{urg\_box\_xyz\_m} & 30,30,8 & rolling local field。\\
  \bottomrule
\end{longtable}
\end{center}

% ============================================================
\section{数据结构与接口规格}
% ============================================================

\subsection{Snapshot 输入}

Advisory predictor 输入应包含：
\begin{itemize}
\item current pose $\hat{T}_t$；
\item position covariance $\Sig^p_t$ 或 pose covariance $\Sig^{pose}_{6\times6}$；
\item visible satellite epoch：satellite position、elevation、azimuth、variance；
\item local ESDF/occupancy map；
\item LiDAR map voxel normal statistics；
\item current LiDAR ARAIM block quality summary only for bias proxy/logging，不用于 certified claim。
\end{itemize}


\subsection{Fallback rules for missing or invalid inputs}

\begin{center}
\small
\begin{tabularx}{\linewidth}{P{0.20\linewidth}Y Y}
  \toprule
  \textbf{Input/source} & \textbf{Failure mode} & \textbf{Required fallback} \\
  \midrule
  Prior covariance & missing, NaN/Inf, non-positive eigenvalue & use configured diagonal prior covariance; clamp eigenvalues; set \texttt{FIM\_REGULARIZED}; log once per refresh \\
  GNSS satellites & satellite position unavailable & skip satellite; if remaining satellites $<4$, set $\Lam^G_{pos}=0$ and \texttt{GNSS\_DEGENERATE} \\
  GNSS variance & missing or non-positive variance & use \texttt{sigma\_gnss\_default\_m}; set warning counter in refresh log \\
  LiDAR normals & voxel normal statistics unavailable & use old advisory scalar predictor if configured, otherwise $\Lam^L_{pos}=0$ and \texttt{LIDAR\_DEGENERATE} \\
  LiDAR neighbors & fewer than \texttt{lidar\_min\_valid\_normals} valid normals & set \texttt{LIDAR\_DEGENERATE}; do not invert untrusted LiDAR-only FIM \\
  FIM inverse & condition number too large, NaN/Inf detected & add eigenvalue floor / regularization; set \texttt{FIM\_REGULARIZED}; clip output PL/cost \\
  ESDF/AL & ESDF query invalid or out of range & set \texttt{OUT\_OF\_RANGE}; use conservative low AL or unknown-risk policy according to config \\
  Stale field & age $>t_{max\_age}$ & set \texttt{STALE\_PL} and \texttt{UNKNOWN\_RISK}; apply unknown penalty; trigger refresh/slowdown hook \\
  \bottomrule
\end{tabularx}
\end{center}

\subsection{UnifiedRiskVoxel suggested C++ struct}

\begin{lstlisting}[language=C++]
struct UnifiedRiskVoxel {
  float esdf_m = std::numeric_limits<float>::quiet_NaN();
  float occ_prob = 0.5f;

  float al_m = 0.0f;
  float hpl_adv_m = 0.0f;
  float vpl_adv_m = 0.0f;
  float pl_adv_m = 0.0f;
  float integrity_margin_m = 0.0f;
  float pi_cost = 0.0f;

  float age_s = 0.0f;
  uint32_t flags = 0u;
};
\end{lstlisting}

\subsection{Flags}

\begin{lstlisting}[language=C++]
enum UnifiedRiskFlags : uint32_t {
  VALID_ESDF       = 1u << 0,
  VALID_OCCUPANCY  = 1u << 1,
  VALID_PL         = 1u << 2,
  STALE_PL         = 1u << 3,
  UNKNOWN_RISK     = 1u << 4,
  OCCUPIED         = 1u << 5,
  OUT_OF_RANGE     = 1u << 6,
  GNSS_DEGENERATE  = 1u << 7,
  LIDAR_DEGENERATE = 1u << 8,
  FIM_REGULARIZED      = 1u << 9,
  SS_VARIANCE_FALLBACK = 1u << 10,
  FRAME_MISMATCH       = 1u << 11,
  CLOCK_UNIT_FALLBACK  = 1u << 12
};
\end{lstlisting}

\subsection{Query API}

\begin{lstlisting}[language=C++]
struct RiskQueryResult {
  bool valid = false;
  Eigen::Vector3d grad_pi = Eigen::Vector3d::Zero();
  double esdf_m = 0.0;
  double al_m = 0.0;
  double hpl_adv_m = 0.0;
  double vpl_adv_m = 0.0;
  double pl_adv_m = 0.0;
  double integrity_margin_m = 0.0;
  double pi_cost = 0.0;
  double age_s = 0.0;
  uint32_t flags = 0u;
};

RiskQueryResult queryRisk(const Eigen::Vector3d& p_W,
                          bool need_gradient) const;
\end{lstlisting}

\subsection{Topic/message naming guideline}

Codex 应 inspect 当前 message definitions 后复用最小改动。若新增字段，推荐命名：
\begin{center}
\small
\begin{tabularx}{\linewidth}{lY}
  \toprule
  \textbf{Name} & \textbf{Meaning} \\
  \midrule
  \texttt{/iap/integrity/current\_pl} & current certified monitor PL and alarm。\\
  \texttt{/iap/advisory/pl\_field} & advisory predicted PL field。\\
  \texttt{/iap/advisory/risk\_grid} & URG multi-layer risk grid。\\
  \texttt{/iap/planner/integrity\_cost} & planner-consumable PI cost samples if URG not yet integrated。\\
  \texttt{hpl\_certified\_m}, \texttt{vpl\_certified\_m} & monitor output。\\
  \texttt{hpl\_adv\_m}, \texttt{vpl\_adv\_m}, \texttt{pl\_adv\_m} & advisory output。\\
  \texttt{integrity\_margin\_m} & $\AL-\PL^{adv}$。\\
  \bottomrule
\end{tabularx}
\end{center}


% ============================================================
\section{跨阶段不变项 / Stage Invariants}
\label{sec:stage-invariants-placeholder}
% ============================================================

以下不变项优先级高于任何局部实现建议。若 Codex 发现实际代码结构与本 spec 冲突，必须停止并报告，不得自行解释后继续修改。

\begin{center}
\small
\begin{tabularx}{\linewidth}{P{0.16\linewidth}Y Y}
  \toprule
  \textbf{Stage} & \textbf{Invariant} & \textbf{Reason} \\
  \midrule
  All & Current certified monitor fusion remains $\PL^{mon}_q=\max(\PL^G_q,\PL^L_q)$. & Preserves conservative alarm path and avoids unmodeled joint fault tree. \\
  All & Advisory predictor output must be named/logged as \texttt{adv}, \texttt{proxy}, or \texttt{advisory}; never as certified PL. & Prevents certification overclaim. \\
  Stage 0 & Planner behavior, PI weights, and cost field topics must remain unchanged. & Stage 0 isolates LiDAR PL stability. \\
  Stage 1 & Numerical behavior should not change except logging/naming. & Stage 1 is semantic cleanup only. \\
  Stage 2 & FIM-add is implemented only in advisory predictor; old predictor mode remains available for A/B test. & Enables regression comparison and preserves monitor path. \\
  Stage 3 & ESDF/collision/dynamic weights must not be tuned automatically. & Prevents PI benefits from being confounded with planner retuning. \\
  Stage 4 & Old \texttt{/iap/integrity\_front\_cost\_field} and \texttt{/iap/integrity\_cost\_field} remain published until URG migration is validated. & Maintains launch/demo compatibility. \\
  Stage 5 & Performance optimization cannot change numerical semantics without A/B logs. & Prevents cache/parallel bugs from changing planner behavior silently. \\
  \bottomrule
\end{tabularx}
\end{center}

\begin{warningbox}{Hard stop condition}
  If any stage requires changing the current monitor fusion rule, deleting compatibility topics, changing frame conventions, or changing collision cost weights, Codex must stop and ask for approval.
\end{warningbox}

% ============================================================
\section{阶段化实现计划}
% ============================================================

\subsection{Stage 0：LiDAR ARAIM PL stability repair}

\subsubsection{Goal}
消除 LiDAR certified PL 在静止/匀速场景下的非物理线性增长，同时不改变 GNSS ARAIM、不改变 planner 行为。

\subsubsection{Tasks}
\begin{enumerate}
\item Locate LiDAR block risk score computation and bias overbound computation.
\item Replace unbounded age term by bounded model.
\item Enforce fixed target keyframe sliding window $|T_t|\le K_T$.
\item Add PL component logging: $|d_f|$, $K_{fa}\sigma_{ss}$, $K_{md}\sigma_f$, $b_f$.
\item Add gamma component logging: RMSE, inlier, condition, age.
\item Add config parameters with defaults from Section 4.
\item Add sanity check script or log command for 10 min static run.
\end{enumerate}

\subsubsection{Do not}
\begin{itemize}
\item Do not change certified monitor fusion.
\item Do not introduce FIM-add.
\item Do not change planner cost.
\item Do not rename topics in this stage unless unavoidable.
\end{itemize}

\subsubsection{Acceptance}
\begin{acceptbox}{Stage 0 acceptance}
\begin{itemize}
  \item Static 10 min run: LiDAR PL RMS drift $<0.5$ m or configured threshold.
  \item No monotonic unbounded linear trend in $b_f$ or total PL.
  \item Logs show PL components and gamma components separately.
  \item Current alarm behavior remains compatible with pre-change baseline except removal of artificial drift.
\end{itemize}
\end{acceptbox}

\subsection{Stage 1：Naming and interface clarification}

\subsubsection{Goal}
只做语义正名，不改变核心算法行为。区分 current certified outputs 和 future advisory outputs。

\subsubsection{Tasks}
\begin{enumerate}
\item Rename comments/log labels from generic ``Predicted PL'' to ``Advisory PL Proxy'' where applicable.
\item Add documentation comments to current monitor outputs: ``certified/current monitor path''.
\item Add documentation comments to predictor outputs: ``advisory/future planning path; non-certified''.
\item Ensure planner consumes advisory fields only, not monitor alarm outputs directly.
\end{enumerate}

\subsubsection{Acceptance}
\begin{acceptbox}{Stage 1 acceptance}
\begin{itemize}
  \item Code and logs no longer confuse current certified PL and advisory PL proxy.
  \item Existing launch files still run.
  \item No numerical behavior change expected except names/logs/comments.
\end{itemize}
\end{acceptbox}

\subsection{Stage 2：Advisory FIM predictor and FIM-add fusion}

\subsubsection{Goal}
将 future predictor 从 two-channel PL proxy/max logic 升级为 position-FIM-level fusion：
\[
\Lam^{\adv}_{pos}=\Lam^{\prior}_{pos}+\Lam^G_{pos}+\Lam^L_{pos}.
\]

\subsubsection{Stage 2.1：GNSS position FIM}
\begin{enumerate}
\item Build $G\in\mathbb{R}^{N\times4}$ from visible satellites at candidate position.
\item Build $W=\diag(\sigma^{-2}_{eff,s})$ using existing variance model.
\item Compute $H_G=G^\top W G$.
\item Convert to position FIM using Schur complement:
  \[
    \Lam^G_{pos}=H_{pp}-H_{pc}(H_{cc}+\lambda_c)^{-1}H_{cp}.
  \]
\item Handle degeneracy with flags and zero contribution.
\end{enumerate}

\subsubsection{Stage 2.2：LiDAR \texorpdfstring{$3\times3$}{3x3} FIM predictor}
\begin{enumerate}
\item Extend map voxel data to store normal histogram/statistics, or build a side cache.
\item For query point $\pvec$, collect neighboring voxels within $R_L$.
\item Compute $\pi_m(\pvec)$ and $w_m(\pvec)$.
\item Accumulate $\Lam^L_{pos}=\sum \pi_m w_m nn^\top$.
\item Detect degeneracy by eigenvalues and condition number:
  \[
    \lambda_{min}(\Lam^L_{pos}) < \lambda_L^{min}
    \quad \Rightarrow \quad
    \texttt{LIDAR\_DEGENERATE}.
  \]
\end{enumerate}

\subsubsection{Stage 2.3：FIM-add fusion and PL extraction}
\begin{enumerate}
\item Build prior FIM from covariance.
\item Add prior, GNSS, LiDAR FIM.
\item Invert with regularization.
\item Compute HPL/VPL/advisory scalar PL.
\item Preserve old predictor mode under \texttt{fusion\_mode=max} for A/B test.
\end{enumerate}

\subsubsection{Do not}
\begin{itemize}
\item Do not modify certified/current ARAIM monitor fusion or alarm logic.
\item Do not remove the old predictor mode; keep \texttt{fusion\_mode=max} or equivalent for A/B tests.
\item Do not require raycast in the first LiDAR FIM implementation.
\item Do not introduce full $6\times6$ LiDAR pose FIM in Stage 2.
\item Do not change planner cost or planner topic consumption in Stage 2.
\end{itemize}

\subsubsection{Acceptance}
\begin{acceptbox}{Stage 2 acceptance}
\begin{itemize}
  \item Corridor scenario: along-corridor PL significantly larger than cross-corridor PL, anisotropy ratio $>3$.
  \item Open-sky GNSS-rich region: GNSS FIM improves advisory PL.
  \item LiDAR-feature-rich region: LiDAR FIM improves advisory PL.
  \item Degenerate cases produce flags, not NaNs.
  \item A/B switch \texttt{fusion\_mode=max/fim\_add} works.
\end{itemize}
\end{acceptbox}

\subsection{Stage 3：PI function redesign}

\subsubsection{Goal}
PI function 只消费统一 advisory PL/AL/IM，不再直接处理 multi-source GNSS/LiDAR PL channels。

\subsubsection{Tasks}
\begin{enumerate}
\item Implement base cost:
  \[
    c_{PI}^{main}=\lambda_\pi[\PL^{adv}-\AL+m]^2_+.
  \]
\item Add optional ratio term behind config flag.
\item Add optional balance/fault-robust terms behind config flags; default off.
\item Clip cost to \texttt{pi\_max\_cost}.
\item Ensure A* front-end and B-spline back-end query same PI interface.
\item Add logs: AL, PLadv, IM, cPI, active terms.
\end{enumerate}

\subsubsection{Do not}
\begin{itemize}
\item Do not enable ratio, balance, or fault-robust terms by default.
\item Do not silently tune planner weights or collision/ESDF weights.
\item Do not consume current certified monitor PL as the planner field.
\item Do not change A* and B-spline backend behavior in separate ways; both must use the same PI interface after Stage 3.
\end{itemize}

\subsubsection{Acceptance}
\begin{acceptbox}{Stage 3 acceptance}
\begin{itemize}
  \item Planner works with $c_{PI}^{main}$ alone.
  \item In a low-IM region, A* initial path is biased away from risk when feasible.
  \item Turning off PI returns standard EGO behavior.
  \item Balance/fault-robust terms do not affect baseline when disabled.
\end{itemize}
\end{acceptbox}

\subsection{Stage 4：Unified Risk Grid}

\subsubsection{Goal}
合并 ESDF/occupancy/integrity field，避免重复查询与 stale fallback 语义错误。Stage 4 必须拆成小 patch，不得一次性重构 planner。

\subsubsection{Stage 4.1：URG data structure and publisher compatibility}
\begin{enumerate}
\item Implement \texttt{UnifiedRiskVoxel} or equivalent minimal extension.
\item Add URG rolling grid bounds/resolution configs.
\item Fill ESDF, occupancy, AL, PLadv, IM, PI cost, age, flags.
\item Keep old planner-consumable topics alive through compatibility publishing.
\end{enumerate}

\subsubsection{Stage 4.2：single query API}
\begin{enumerate}
\item Implement a single interpolator / \texttt{queryRisk()} API.
\item Validate query output against existing ESDF and PI field queries.
\item Do not switch both A* and B-spline backend in the same patch.
\end{enumerate}

\subsubsection{Stage 4.3：planner integration}
\begin{enumerate}
\item First migrate A* front-end to \texttt{queryRisk()}.
\item Then migrate B-spline backend to \texttt{queryRisk()}.
\item Replace stale-to-zero fallback with unknown penalty or refresh/slowdown flag.
\end{enumerate}

\subsubsection{Do not}
\begin{itemize}
\item Do not delete old topics until URG output is validated.
\item Do not replace A* and B-spline query paths in one patch.
\item Do not let stale field return zero PI cost silently.
\end{itemize}

\subsubsection{Acceptance}
\begin{acceptbox}{Stage 4 acceptance}
\begin{itemize}
  \item A* and back-end read from the same query API after staged migration.
  \item Old topics remain available during migration.
  \item Stale field never becomes zero risk silently.
  \item Debug visualization can show PL, AL, IM, PI cost, stale flags.
  \item Query latency decreases or remains no worse than old two-grid approach.
\end{itemize}
\end{acceptbox}

\subsection{Stage 5：Performance optimization and regression}

\subsubsection{Goal}
在功能稳定后进行性能优化，目标 PL/URG refresh $<100$ ms。

\subsubsection{Allowed optimizations}
\begin{itemize}
\item sparse active voxels only；
\item incremental LiDAR FIM accumulation；
\item GNSS FIM cache by satellite geometry / local candidate batch；
\item multi-resolution grid；
\item parallel grid refresh；
\item Eigen fixed-size matrix optimization for $3\times3$ and $4\times4$ operations。
\end{itemize}

\subsubsection{Acceptance}
\begin{acceptbox}{Stage 5 acceptance}
\begin{itemize}
  \item Local URG refresh $<100$ ms for default $30\times30\times8$ m, 1 m resolution.
  \item Planner loop remains stable at target rate.
  \item No NaN/Inf in FIM, PL, PI cost over long run.
  \item Full regression experiments pass.
\end{itemize}
\end{acceptbox}

% ============================================================
\section{Suggested commit plan}
% ============================================================

\begin{center}
\small
\begin{tabularx}{\linewidth}{lY Y}
  \toprule
  \textbf{Commit} & \textbf{Scope} & \textbf{Expected output} \\
  \midrule
  \texttt{stage0-lidar-pl-stability} & bounded age, fixed target window, PL component logs & static PL drift removed。\\
  \texttt{stage1-advisory-naming} & comments/topic/log naming cleanup & certified/advisory distinction clear。\\
  \texttt{stage2a-gnss-pos-fim} & GNSS $4\times4$ FIM + Schur complement & $\Lam^G_{pos}$ query works。\\
  \texttt{stage2b-lidar-trans-fim} & voxel normal FIM predictor & corridor anisotropy appears。\\
  \texttt{stage2c-fim-add-pl} & FIM fusion + PLadv extraction & fim\_add mode works。\\
  \texttt{stage3-pi-main-cost} & unified PI function & planner consumes PLadv only。\\
  \texttt{stage3-ablation-terms} & optional balance/fault-robust, default off & ablation switches work。\\
  \texttt{stage4a-urg-compat-grid} & URG struct/fill + old topic compatibility & risk grid available without breaking demo。\\
  \texttt{stage4b-risk-query-api} & single interpolator/queryRisk API & A/B query consistency checked。\\
  \texttt{stage4c-planner-query-migration} & migrate A* then B-spline to queryRisk & single query API in planner。\\
  \texttt{stage5-perf-regression} & profiling/cache/parallelization & refresh $<100$ ms。\\
  \bottomrule
\end{tabularx}
\end{center}

% ============================================================
\section{Logging and debugging requirements}
% ============================================================

\subsection{Required logs for Stage 0}

Every current LiDAR ARAIM epoch should optionally log:
\begin{lstlisting}
time_s, hpl_lidar, vpl_lidar,
worst_hypothesis_type, worst_hypothesis_id,
sep_term_m, sigma_ss_term_m, sigma_subset_term_m, bias_term_m,
gamma_rmse, gamma_inlier, gamma_condition, gamma_age,
target_window_size, lambda_min_subset, condition_number_subset,
sigma_ss_raw_m2, sigma_ss_used_m, ss_variance_fallback_flag
\end{lstlisting}

\subsection{Required logs for Stage 2/3}

Every advisory field refresh should log:
\begin{lstlisting}
time_s, num_grid_points, refresh_ms,
num_gnss_visible_mean, num_lidar_voxels_mean,
pl_adv_min, pl_adv_mean, pl_adv_p95, pl_adv_max,
al_min, al_mean, im_min, im_mean,
num_gnss_degenerate, num_lidar_degenerate, num_fim_regularized,
num_sat_skipped, num_missing_variance, num_lidar_fallback,
frame_map_name, frame_gnss_local_name, frame_ecef_to_local_available, gnss_clock_unit,
mode_fusion, mode_lidar_fim, use_ratio, use_balance, use_fault_robust
\end{lstlisting}

\subsection{Required debug visualizations}

\begin{itemize}
\item PL field heatmap.
\item AL field heatmap.
\item IM field heatmap.
\item PI cost field heatmap.
\item GNSS-only PL vs LiDAR-only PL vs fused PL.
\item LiDAR FIM eigenvectors/eigenvalues in corridor scenario.
\item Stale/unknown risk flags.
\end{itemize}

% ============================================================
\section{实验与验收矩阵}
% ============================================================

\begin{center}
\scriptsize
\begin{longtable}{@{}P{0.05\linewidth}P{0.17\linewidth}P{0.22\linewidth}P{0.34\linewidth}P{0.05\linewidth}@{}}
  \toprule
  \textbf{ID} & \textbf{Experiment} & \textbf{Purpose} & \textbf{Acceptance} & \textbf{Stage} \\
  \midrule
  A1 & Static LiDAR PL stability & Detect artificial PL drift & 10 min static RMS drift $<0.5$ m, no linear trend & 0 \\
  A2 & Constant-velocity PL stability & Detect motion-related drift & PL bounded, no monotonic age-driven growth & 0 \\
  A3 & GNSS fault injection & Certified monitor still works & GNSS PL rises, max monitor alarm unchanged & 0/2 \\
  A4 & LiDAR block fault injection & LiDAR monitor still works & target/source/level hypothesis produces higher PL & 0 \\
  B1 & Predictor self-consistency at current pose & Advisory field sanity & $\PL^{adv}(p_t)$ close to calibrated monitor PL proxy & 2 \\
  B2 & Corridor anisotropy & LiDAR FIM geometry & Along-corridor PL $>$ cross-corridor PL, ratio $>3$ & 2 \\
  B3 & Open sky vs canopy/canyon & GNSS geometry & Better sky view gives lower GNSS advisory PL & 2 \\
  B4 & FIM-add vs max A/B & Fusion benefit & FIM-add tighter but no NaN/false unsafe collapse & 2 \\
  C1 & PI main cost only & Planner integration & Avoids low-IM area when feasible; off returns baseline & 3 \\
  C2 & Ratio term sensitivity & Tune optional term & No instability; ratio term optional & 3 \\
  C3 & Fault-robust ablation & Source-redundancy value & Helps under forced single-source loss; default off & 3 \\
  D1 & URG stale field test & Unknown risk semantics & Frozen field triggers unknown penalty/refresh, not zero risk & 4 \\
  D2 & Query latency & Runtime performance & URG query faster/no worse than old dual lookup & 4 \\
  E1 & Full regression & End-to-end stability & No crash/NaN; refresh $<100$ ms after optimization & 5 \\
  \bottomrule
\end{longtable}
\end{center}


% ============================================================
\section{Codex execution gate：每个 stage 的强制流程}
% ============================================================

Codex 不得从本文档直接实现 Stage 2/3/4。每个 stage 都必须先经过以下 gate。

\begin{enumerate}
  \item \textbf{Inspect only:} 输出实际 repository 中的文件、类、函数、topic、message、config、launch/test/log 命令。
  \item \textbf{Frame audit:} 明确 planner frame、ESDF/map frame、LiDAR map frame、GNSS local frame、ECEF/ENU/map transform 工具。
  \item \textbf{Patch boundary:} 说明本 stage 会改哪些文件，不会改哪些文件。
  \item \textbf{Invariant check:} 逐条确认 Section~\ref{sec:stage-invariants-placeholder} 的不变项未被破坏。若 LaTeX label 不存在，仍需引用 ``Stage Invariants'' section。
  \item \textbf{Validation plan:} 给出 build/test/launch/log/plot 命令。没有可运行命令时，必须说明缺口。
  \item \textbf{Approval required:} 未得到人工批准，不得修改代码。
\end{enumerate}

\begin{codexbox}{Generic stage-gate prompt}
\begin{lstlisting}
Before implementing this stage, inspect the repository and produce a file-level sub-plan.
Do not edit code yet. Your sub-plan must include:
1. exact files/classes/functions to modify,
2. existing topics/messages/configs to preserve,
3. build/test/launch/log commands,
4. coordinate-frame assumptions and transform utilities,
5. stage invariants that must remain unchanged,
6. fallback behavior and debug logs,
7. risks and unknowns.
Wait for approval before making any code changes.
\end{lstlisting}
\end{codexbox}

% ============================================================
\section{Initial Codex prompt}
% ============================================================

\begin{codexbox}{Prompt to provide to Codex}
\begin{lstlisting}
You are working on the IAP repository. Read this implementation plan first.
Your first task is NOT to modify code. First inspect the repository and produce
an exact file-level implementation plan for Stage 0 only.

Hard constraints:
1. Do not change the certified/current ARAIM monitor fusion rule:
   PL_mon = max(PL_G, PL_L).
2. Do not introduce FIM-add into the certified monitor path.
3. Stage 0 only: fix LiDAR PL stability by bounded age, fixed target window,
   and component logging.
4. Do not modify planner behavior in Stage 0.
5. Identify the exact files/classes/functions where LiDAR block risk score,
   target keyframe selection, PL composition, and logging are implemented.
6. Identify coordinate frames used by planner, ESDF/map, LiDAR map, and GNSS.
7. Discover the exact build/test/launch/log/plot commands available in this repo.
8. Propose a minimal patch plan and tests/log commands.

After inspection, output:
- relevant files/functions,
- proposed changes,
- config parameters to add,
- logs to add,
- tests or launch commands to run,
- coordinate-frame assumptions,
- fallback/compatibility plan,
- risks/unknowns.
Do not edit code until I approve the plan.
\end{lstlisting}
\end{codexbox}

% ============================================================
\section{Stage-specific Codex prompts}
% ============================================================

\subsection{Stage 0 implementation prompt}
\begin{lstlisting}
Implement Stage 0 only.
- Replace unbounded LiDAR age risk with bounded age model.
- Add fixed target keyframe sliding window K.
- Add PL component logs and gamma component logs.
- Add config parameters with safe defaults.
- Do not change certified GNSS/LiDAR fusion.
- Do not change planner behavior.
- Build and run available tests or provide exact commands if runtime simulation is needed.
Return changed files, summary, and validation result.
\end{lstlisting}

\subsection{Stage 2 implementation prompt}
\begin{lstlisting}
Implement Stage 2 advisory FIM fusion only after Stage 0/1 are merged.
- Add GNSS position FIM using G^T W G and Schur complement to remove clock.
- Add LiDAR 3x3 translational FIM from voxel normal statistics.
- Add FIM-add fusion in advisory predictor only.
- Derive HPL_adv/VPL_adv/PL_adv from position covariance.
- Keep old max mode behind fusion_mode=max for A/B tests.
- Degenerate FIM must set flags and fallback safely; never produce NaN/Inf.
- If LiDAR normal statistics are unavailable, use configured fallback instead of refactoring the map.
- Log coordinate-frame assumptions for GNSS, LiDAR, ESDF, and planner queries.
- Do not change planner cost or certified/current ARAIM monitor fusion.
\end{lstlisting}

\subsection{Stage 3/4 implementation prompt}
\begin{lstlisting}
Implement Stage 3 or one Stage 4 sub-stage only after Stage 2 passes corridor anisotropy test.
Stage 3:
- Redesign PI function to consume unified PL_adv and AL only.
- Base PI cost is hinge margin cost; ratio/balance/fault-robust default off.
- Do not change ESDF/collision cost or certified monitor outputs.
Stage 4:
- First build Unified Risk Grid with ESDF, occupancy, AL, HPL_adv, VPL_adv,
  PL_adv, IM, PI cost, age, flags, while preserving old topics.
- Then implement queryRisk API and validate against old query outputs.
- Migrate A* first, then B-spline backend; do not migrate both in one patch.
- Replace stale-to-zero behavior with unknown-risk penalty or refresh/slowdown flag.
\end{lstlisting}

% ============================================================
\section{Known risks and mitigation}
% ============================================================

\begin{center}
\small
\begin{tabularx}{\linewidth}{lY Y}
  \toprule
  \textbf{Risk} & \textbf{Failure mode} & \textbf{Mitigation} \\
  \midrule
  GNSS/LiDAR FIM scale mismatch & One source dominates artificially & Use covariance inflation, log source contributions, keep ablation off by default。\\
  LiDAR FIM degeneracy & Matrix inversion NaN/Inf & Regularize FIM, eigenvalue floor, flags。\\
  Predicted PL overclaim & Reviewer questions certification & Always call future output advisory PL proxy; certified monitor remains separate。\\
  PI over-conservatism & Planner refuses feasible routes & Start with hinge main cost only, tune weights, keep ratio/fault terms optional。\\
  Stale field semantics & Old risk treated as safe & Unknown penalty and stale flags。\\
  Coordinate-frame mismatch & PL field appears shifted or GNSS/LiDAR FIM wrong & Log frames, reuse existing transforms, stop if transform missing。\\
  Missing source data & Crash or NaN when satellite/normal/covariance unavailable & Explicit fallback rules and flags。\\
  Interface breakage & Existing demo topics disappear & Compatibility publisher and old topic retention until Stage 4 validation。\\
  Codex over-refactor & Breaks running demo & Stage-wise commits, inspect-first workflow, no cross-stage edits。\\
  \bottomrule
\end{tabularx}
\end{center}

% ============================================================
\section{Final success criteria}
% ============================================================

本整改完成后，系统应满足：
\begin{enumerate}
\item Current certified monitor path 仍保持 conservative max fusion，alarm 行为可解释。
\item LiDAR certified PL 不再因 age 或 growing target set 出现无界线性增长。
\item Future predictor 输出清晰标注为 advisory PL proxy，而不是 certified PL。
\item Advisory predictor 使用 position-domain FIM-add 融合 GNSS、LiDAR、prior。
\item LiDAR predictor 能在走廊中表现出方向性不确定性。
\item PI function 只消费统一 $\PL^{adv}$、$\AL$、$\IM$，而不是多源 PL 拼接。
\item URG 提供单一查询接口给 A* 和 B-spline optimizer。
\item Stale PL field 不再 fallback 到 zero risk。
\item 坐标系假设在代码注释、日志和 Codex plan 中明确记录。
\item 旧 cost-field topics 在 URG 完成前保持兼容。
\item 默认配置可跑通 demo；ablation 项默认关闭。
\item PL field / URG refresh 经优化后达到 $<100$ ms 目标。
\end{enumerate}

\vspace{1em}
\begin{center}
\Large
\fbox{\parbox{0.9\linewidth}{\centering
\textbf{Core message:}\quad Current monitoring remains conservative; future planning uses fused advisory information.}}
\end{center}

\end{document}
